% GENERATED FILE. Edit canonical lessons, then run npm run generate:books.
% canonical-source-hash: fnv1a64:ffb9e3b4b5d1fde1
% canonical-lessons: TA-C158-camipattil, TA-C158-erkanave, TA-C158-pona, TA-C158-oru-kalattil, TA-C158-kali

\chapter{Recently, Already, Last, Once, To Spend (Time)}
\label{ch:ta-recently-already-last-once-to-spend-time}
\hlchaptermodality{158}

\begin{chapteropening}
\textbf{By the end of this chapter:} \emph{I can say what happened, and when: recently, already, last week.}

Complete the last lesson of chapter 158: I can say what happened, and when: recently, already, last week.
\end{chapteropening}

\section[\texorpdfstring{\ta{சமீபத்தில்} (camīpattil)}{சமீபத்தில் (camīpattil)}]{\ta{சமீபத்தில்} (camīpattil) — recently}

\label{lesson:TA-C158-camipattil}



\begin{quote}
Before the new one: say the Tamil for to burn, then the Tamil for to put out.
\end{quote}

\subsection*{You'll want to know: \ta{சமீபத்தில்}}
\textbf{\ta{சமீபத்தில்}} — \emph{camīpattil} — \textquotedblleft{}recently\textquotedblright{}. Tamil's past for \textquotedblleft{}I\textquotedblright{} ends in \textbf{-\ta{ஏன்}} after a past stem: \textbf{\ta{நான்} \ta{சமீபத்தில்} \ta{ஒரு} \ta{புத்தகம்} \ta{படித்தேன்}} — \emph{nāṉ camīpattil oru puttakam paṭittēṉ} — \textquotedblleft{}I recently read a book.\textquotedblright{}

\begin{grammarlens}[title={What you've built}]
One more word for this chapter.
\end{grammarlens}

\subsection*{Your turn}
\begin{itemize}
\raggedright
  \item \emph{Say it:} \emph{camīpattil}
  \item \emph{Say it:} \emph{camīpattil}, once more
  \item \emph{Say it:} say \emph{camīpattil}
  \item \emph{Recall:} say the Tamil for to burn, then the Tamil for to put out, then say \emph{camīpattil} again
\end{itemize}

\begin{tcolorbox}[breakable,colback=teal!4,colframe=teal!35!black,title={Before you move on}]
What does \ta{சமீபத்தில்} mean? (\textquotedblleft{}Recently\textquotedblright{}.) Say it once more.
\end{tcolorbox}
\section[\texorpdfstring{\ta{ஏற்கனவே} (ēṟkaṉavē)}{ஏற்கனவே (ēṟkaṉavē)}]{\ta{ஏற்கனவே} (ēṟkaṉavē) — already}

\label{lesson:TA-C158-erkanave}



\begin{quote}
Before the new one: say the Tamil for to put out, then the Tamil for recently.
\end{quote}

\subsection*{You'll want to know: \ta{ஏற்கனவே}}
\textbf{\ta{ஏற்கனவே}} — \emph{ēṟkaṉavē} — \textquotedblleft{}already\textquotedblright{}. \textbf{\ta{நான்} \ta{ஏற்கனவே} \ta{சாப்பிட்டேன்}} — \emph{nāṉ ēṟkaṉavē cāppiṭṭēṉ} — \textquotedblleft{}I have already eaten.\textquotedblright{}

\begin{grammarlens}[title={What you've built}]
One more word for this chapter.
\end{grammarlens}

\subsection*{Your turn}
\begin{itemize}
\raggedright
  \item \emph{Say it:} \emph{ēṟkaṉavē}
  \item \emph{Say it:} \emph{ēṟkaṉavē}, once more
  \item \emph{Say it:} say \emph{ēṟkaṉavē}
  \item \emph{Recall:} say the Tamil for to put out, then the Tamil for recently, then say \emph{ēṟkaṉavē} again
\end{itemize}

\begin{tcolorbox}[breakable,colback=teal!4,colframe=teal!35!black,title={Before you move on}]
What does \ta{ஏற்கனவே} mean? (\textquotedblleft{}Already\textquotedblright{}.) Say it once more.
\end{tcolorbox}
\section[\texorpdfstring{\ta{போன} (pōṉa)}{போன (pōṉa)}]{\ta{போன} (pōṉa) — last, previous}

\label{lesson:TA-C158-pona}



\begin{quote}
Before the new one: say the Tamil for recently, then the Tamil for already.
\end{quote}

\subsection*{You'll want to know: \ta{போன}}
\textbf{\ta{போன}} — \emph{pōṉa} — \textquotedblleft{}last, previous\textquotedblright{}. It is the past form of \textbf{\ta{போ}}, \textquotedblleft{}to go\textquotedblright{}, used before a noun: \textbf{\ta{போன} \ta{வாரம்}} — \emph{pōṉa vāram} — \textquotedblleft{}last week\textquotedblright{}; \textbf{\ta{போன} \ta{வருடம்}} — \emph{pōṉa varuṭam} — \textquotedblleft{}last year\textquotedblright{}.

\begin{grammarlens}[title={What you've built}]
One more word for this chapter.
\end{grammarlens}

\subsection*{Your turn}
\begin{itemize}
\raggedright
  \item \emph{Say it:} \emph{pōṉa}
  \item \emph{Say it:} \emph{pōṉa}, once more
  \item \emph{Say it:} say \emph{pōṉa}
  \item \emph{Recall:} say the Tamil for recently, then the Tamil for already, then say \emph{pōṉa} again
\end{itemize}

\begin{tcolorbox}[breakable,colback=teal!4,colframe=teal!35!black,title={Before you move on}]
What does \ta{போன} mean? (\textquotedblleft{}Last, previous\textquotedblright{}.) Say it once more.
\end{tcolorbox}
\section[\texorpdfstring{\ta{ஒரு} \ta{காலத்தில்} (oru kālattil)}{ஒரு காலத்தில் (oru kālattil)}]{\ta{ஒரு} \ta{காலத்தில்} (oru kālattil) — once, at one time}

\label{lesson:TA-C158-oru-kalattil}



\begin{quote}
Before the new one: say the Tamil for already, then the Tamil for last.
\end{quote}

\subsection*{You'll want to know: \ta{ஒரு} \ta{காலத்தில்}}
\textbf{\ta{ஒரு} \ta{காலத்தில்}} — \emph{oru kālattil} — \textquotedblleft{}once, at one time\textquotedblright{}. It points far back: \textbf{\ta{ஒரு} \ta{காலத்தில்} \ta{இங்கே} \ta{ஒரு} \ta{குளம்} \ta{இருந்தது}} — \emph{oru kālattil iṅkē oru kuḷam iruntatu} — \textquotedblleft{}There was once a pond here.\textquotedblright{} \textbf{\ta{இருந்தது}} is the past of \textbf{\ta{இரு}}.

\begin{grammarlens}[title={What you've built}]
One more word for this chapter.
\end{grammarlens}

\subsection*{Your turn}
\begin{itemize}
\raggedright
  \item \emph{Say it:} \emph{oru kālattil}
  \item \emph{Say it:} \emph{oru kālattil}, once more
  \item \emph{Say it:} say \emph{oru kālattil}
  \item \emph{Recall:} say the Tamil for already, then the Tamil for last, then say \emph{oru kālattil} again
\end{itemize}

\begin{tcolorbox}[breakable,colback=teal!4,colframe=teal!35!black,title={Before you move on}]
What does \ta{ஒரு} \ta{காலத்தில்} mean? (\textquotedblleft{}Once, at one time\textquotedblright{}.) Say it once more.
\end{tcolorbox}
\section[\texorpdfstring{\ta{கழி} (kaḻi)}{கழி (kaḻi)}]{\ta{கழி} (kaḻi) — to spend (time)}

\label{lesson:TA-C158-kali}



\begin{quote}
Before the new one: say the Tamil for last, then the Tamil for once.
\end{quote}

\subsection*{You'll want to know: \ta{கழி}}
\textbf{\ta{கழி}} — \emph{kaḻi} — \textquotedblleft{}to spend (time)\textquotedblright{}. Said of time: \textbf{\ta{விடுமுறையைக்} \ta{கிராமத்தில்} \ta{கழித்தேன்}} — \emph{viṭumuṟaiyaik kirāmattil kaḻittēṉ} — \textquotedblleft{}I spent the holiday in the village.\textquotedblright{}

\begin{grammarlens}[title={What you've built}]
One more word for this chapter.
\end{grammarlens}

\subsection*{Your turn}
\begin{itemize}
\raggedright
  \item \emph{Say it:} \emph{kaḻi}
  \item \emph{Say it:} \emph{kaḻi}, once more
  \item \emph{Say it:} say \emph{kaḻi}
  \item \emph{Recall:} say the Tamil for last, then the Tamil for once, then say \emph{kaḻi} again
\end{itemize}

\begin{tcolorbox}[breakable,colback=teal!4,colframe=teal!35!black,title={Before you move on}]
What does \ta{கழி} mean? (\textquotedblleft{}To spend (time)\textquotedblright{}.) Say it once more.
\end{tcolorbox}
